\chapter{From a Hamiltonian to a checked calculation}\label{ch:practical}

\section{State the question and the approximation}

Consider the reference junction of \cref{eq:reference}. We want its ground-state
parity, lowest electron-addition/removal excitation, even pair excitation and
Josephson current. Determining these quantities requires comparing the relevant
sectors and evaluating state-resolved observables. The physical input and finite-bath
approximation are specified in
\code{examples/01_subgap_states/input/reference.json}, shared by the tutorial,
handbook example and figure-generation scripts.

The practical sequence is
\begin{equation*}
 \text{bath and impurity}\ \longrightarrow\ \text{Hamiltonian}\
 \longrightarrow\ \text{sector solves}\ \longrightarrow\ \text{observables and checks}.
\end{equation*}
The bath stores signed normal levels and their contact weights. The Hamiltonian
stores the interacting operator in canonical coordinates. A sector specifies
total quantum numbers of the dot and leads. The result contains state-resolved
energies, vectors or MPS, observables and numerical diagnostics. Physical
parameters, representation choices and solver tolerances have separate roles.

\section{A first calculation}\label{sec:first-calculation}

The following setup is generated from the shared input; it uses explicit
parameters so that it also runs in a fresh Python session. Energies are in gap
units and the QP cutoff is a many-body approximation to the declared finite bath.

% Generated by tools/gen_docs.py from examples/01_subgap_states/input/reference.json.
\begin{lstlisting}
import numpy as np
from qdjj_solver import Sector, cosh_grid, reference_model
from qdjj_solver.qp_solver import SolverOptions, solve

bath = cosh_grid(**{'pairs': 4, 'delta': 1.0, 'bandwidth': 10.0})
model = dict(bath=bath, **{'u': 1.0, 'gamma': 0.5, 'phi': 0.8, 'detuning': 0.2, 'field': 0.0, 'symmetry': False, 'compress': False})
h = reference_model(**model)
even = solve(h, cutoff=2, sector=Sector(0, 0), options=SolverOptions(eigenpairs=2))
odd = solve(h, cutoff=2, sector=Sector(1, 1))
partner = solve(h, cutoff=2, sector=Sector(1, -1))
print(even.energies, odd.energies, partner.energies)
\end{lstlisting}


The last array compares the spin projections of the doublet. Their degeneracy
is a check at zero field; in a finite field both must be included when finding
the global minimum. The even sector supplies the singlet $S$ and excited
singlet $S^*$. At the reference point,
$E_D-E_S=\q{static.doublet_energy}$ and
$E_{S^*}-E_S=\q{static.excited_energy}$. The signed first difference identifies
which parity lies lower. The positive excitation cost is its absolute value,
whereas the pair excitation connects two states of the same even sector.

From the source root the matching even calculation can also be run with:
\begin{lstlisting}[language=bash]
qdjj-solver solve examples/01_subgap_states/input/reference.json --output results/even.json --save-states
\end{lstlisting}
This command solves the sector named by the input. Run the other relevant sectors
to determine the global ground state. The input contains the numerical settings as well as
the model, which makes differences between Python and command-line defaults
irrelevant to this explicit comparison.

\section{Interpret observables and matrix elements}

Each observable array has one entry per returned state. The current derivative
is reported in energy units; multiplying it by $2e/\hbar$ gives the physical
current. For a normalized gap-unit calculation, the corresponding dimensionless
number is $I/(2e\Delta/\hbar)$.

\begin{lstlisting}
current = even.observables["phase_derivative"][0].real
charge = even.observables["impurity_charge"][0].real
addition = odd.matrix_elements("d_up_dagger", even)
pair = even.matrix_elements("current")
print(current, charge, abs(addition[0, 0])**2, abs(pair[1, 0])**2)
\end{lstlisting}

The rows of a transition matrix label the bra/final states and columns the
ket/initial states. Thus the first squared amplitude adds a spin-up dot
electron to $S$ and reaches the chosen odd state, while the second connects
$S$ to $S^*$. The matrix element specifies the coupling strength; predicting
the measured response also requires initial-state populations and line widths. Named electron probes are already
transformed to the model's canonical coordinates. For arbitrary physical
operators use the transformation described in \cref{sec:coordinates}.

\section{Sweep a parameter and retain the results}

The standalone script \code{examples/handbook/scan.py} scans phase, writes both
sector results after every completed point, and retains any failed point with
its diagnostic. Its plotting mode reads those saved results without solving
again. Follow the smooth sector branches first, then construct their lower
envelope and ground-state current as in \cref{fig:primer-current}.

Use a dense phase grid to locate parity crossings, and refine the QP cutoff and
bath to converge their positions. Track states through same-sector crossings
using symmetries, observables and, in compatible coordinates, overlaps.
At exact degeneracy compare the complete relevant subspace.

\section{Ask three separate convergence questions}

\begin{enumerate}
\item Does the result solve its declared finite matrix problem? Inspect residuals
and orthogonality; DMRG also provides sweep and independent-trial diagnostics.
\item Is the many-body representation adequate? Refine QP cutoff or MPS bond
dimension and compare the excitation, current or correlation actually reported.
\item Is the bath adequate? Refine its nodes or fit at fixed physical parameters,
with adequate many-body convergence at each size.
\end{enumerate}
The small DMRG tutorial deliberately begins with insufficient bond dimension,
then approaches the result of full ED for the same finite bath, demonstrating
the first two checks. Apply the third through bath refinement. The larger
reference junction here illustrates a finite-bath, cutoff-two calculation;
assess its continuum accuracy using all three checks.

\tutorial{11_dmrg_convergence}

\section{Save states and prepare for driving}

\begin{lstlisting}
from qdjj_solver import save_result, load_result
save_result(even, "results/even.json", save_states=True)
restored = load_result("results/even.json")
assert np.max(abs(restored.energies - even.energies)) < 1e-12
\end{lstlisting}

The JSON stores scalar results, settings and provenance; the companion NPZ or
HDF5 file stores wavefunctions. Keep both together. Scalar-only plotting can use
ordinary JSON tools, but evaluating a new matrix element requires the states.
Saved, completed DMRG states provide initial wavefunctions for a new compatible
optimization.

The static transition energy and its matrix element now identify a candidate
resonance for the driven part of the manual. Retain fixed canonical coordinates
when introducing the drive. \Cref{ch:floquet,ch:driving} develop that step before
the numerical Floquet methods are applied.

\begin{incode}[The complete software handbook]
\usermanual{docs/static_calculations.md};
\usermanual{docs/workflows.md};
\usermanual{docs/dmrg.md}. The separate User's Guide and Reference Manual PDF
also gives installation, full constructor and solver references, JSON formats,
saved-record fields and troubleshooting. Both manuals are published with the
matching software release.
\end{incode}

\tutorial{00_first_calculation}
